\begin{exercice}\label{exo4-8} (\coolexo)

L'\'etoile la plus brillante du ciel est $\alpha$ du Grand Chien (Sirius). On consid\`ere que la probabilit\'e d'apparition une ann\'ee donn\'ee d'une com\`ete plus brillante que Sirius est \'egale \`a $1/43$. Calculer la probabilit\'e pour que, durant 100 ann\'ees cons\'ecutives, il n'apparaisse aucune com\`ete plus brillante que Sirius ({\it En r\'ealit\'e, les seules com\`etes plus brillantes que Sirius depuis le d\'ebut du 20e si\`ecle ont \'et\'e Halley en 1910 et Hale-Bopp en 1997}).\\
%\corrref{TD1_1}

\end{exercice}
